% Margenes
\geometry{
  a4paper,
  inner=2.5cm,
  outer=2cm,
  top=4cm,
  bottom=3cm,
  headheight=2cm,
  headsep=0.75cm,
  footskip=1.5cm
}

% Espaciado general del texto
% \spacing{1.5}
\setstretch{1.5}

% No estirar espacios verticales para igualar el final de cada página
% (book con twoside usa \flushbottom por defecto → Underfull \vbox)
\raggedbottom

% Eliminar espaciado de enumerate e itemize
\setlist[enumerate,itemize]{itemsep=0pt, topsep=0pt}

% Para que todos los Figura xx aparezcan en negrita
\captionsetup[figure]{labelfont=bf}

% Para que todos los Tabla xx aparezcan en negrita
\DefTblrTemplate{caption-tag}{default}{\bfseries Tabla\hspace{0.25em}\thetable}
\DefTblrTemplate{caption-sep}{default}{\bfseries :\enskip}

% Estilo para páginas de título (como \maketitle)
\fancypagestyle{plain}{
    \fancyhf{}
    \renewcommand{\headrulewidth}{1 pt}
    \renewcommand{\footrulewidth}{1 pt}
    \fancyhead[RO,LE]{Aranda Brian Ezequiel, Wlodeck Fabricio Joaquin}
    \fancyfoot[LE,RO]{\textbf{Pág. N° \thepage}}
}

% Estilo general para el documento
\pagestyle{fancy}
    \fancyhf{}
    \renewcommand{\headrulewidth}{1 pt}
    \renewcommand{\footrulewidth}{1 pt}
    % Encabezado: Siempre a la derecha (R) en todas las páginas
    \fancyhead[LE]{\textbf{\leftmark}}
    \fancyhead[RO]{\textbf{\rightmark}}
    % Pie de página: Número alternado (izquierda en pares, derecha en impares)
    \fancyfoot[LE,RO]{\textbf{Pág. N° \thepage}} 

% Texto parte de las longtable (tablas de varias páginas)
\DefTblrTemplate{contfoot-text}{default}{Continua en la siguiente página $\cdots$}
\DefTblrTemplate{conthead-text}{default}{(Continuación)}

% Definición del espaciado de indice
\setlength{\cftbeforetoctitleskip}{-10 pt}       % Espacio antes de indice general
\setlength{\cftaftertoctitleskip}{5 pt}         % Espacio despues de indice general

% Idem que lo anterior pero para el indice de figuras y tablas
\setlength{\cftbeforeloftitleskip}{-10 pt}       
\setlength{\cftafterloftitleskip}{20 pt}         
\setlength{\cftbeforelottitleskip}{-10 pt}       
\setlength{\cftafterlottitleskip}{20 pt}       

% Cambiar el tamaño de la palabra "Índice"
\renewcommand{\cfttoctitlefont}{\Huge\bfseries} % Tamaño Huge y negrita

% Usar el punto como separador decimal
\decimalpoint

% Configuración de la Bibliografía
\renewcommand{\bibname}{Bibliografía}
\patchcmd{\thebibliography}
  {\chapter*{\bibname}}
  {\chapter{\bibname}\footnotesize}
  {}{}

\begin{comment}
    \titleformat{nombre de la parte a modificar}
        {Formato de la sección}    
        {Número y separador}                             
        {Separación entre numero y tituli}                  
        {Código adicional}
        [Otra configuración, por ejemplo línea]
\end{comment}

% Redefinición del formato de \chapter
\titleformat{\chapter}[hang]
    {\normalfont\huge\bfseries} % sin \centering → alineado a la izquierda
    {}              
    {0 cm}
    {\MakeUppercase}
    \titlespacing*{\chapter}{0 pt}{-30 pt}{5 pt} % {izq}{antes}{después}
    \setcounter{secnumdepth}{-1}
    \setlength{\cftbeforechapskip}{15 pt}

% Redefinición del formato de \section
\titleformat{\section}[block]
    {\normalfont\LARGE\bfseries}    
    {}                    
    {0 cm}                         
    {}                
    %[\vspace{0.5ex}\titlerule] 
    \titlespacing*{\section}{0pt}{0.3 cm}{*1}

% Redefinición del formato de \subsection
\titleformat{\subsection}
    {\normalfont\large\bfseries}    % Formato de la subsección
    {}                              % Elimina el número
    {0 pt}                          % Sangría
    {}                              % Código adicional (opcional)
    \titlespacing*{\subsection}{0pt}{0.2 cm}{*1}

% Redefinición del formato de \subsubsection
\titleformat{\subsubsection}
    {\normalfont\normalsize\bfseries}   % Formato de la subsubsección
    {}                                  % Elimina el número
    {0 pt}                              % Sangría
    {}                                  % Código adicional (opcional)
    \titlespacing*{\subsubsection}{0pt}{0.2 cm}{*1}

% Elimina la numeración del indice
    
    % Secciones
    \renewcommand{\cftsecpresnum}{}         % Elimina el prefijo de numeración
    \renewcommand{\cftsecaftersnum}{}       % Elimina el sufijo de numeración
    \renewcommand{\cftsecnumwidth}{0 pt}    % Ancho de la numeración (0pt para eliminarla)
    \renewcommand{\thesection}{}            % Elimina el número de sección en el índice
    
    % Subsecciones
    \renewcommand{\cftsubsecpresnum}{}          % Elimina el prefijo de numeración
    \renewcommand{\cftsubsecaftersnum}{}        % Elimina el sufijo de numeración
    \renewcommand{\cftsubsecnumwidth}{0 pt}     % Ancho de la numeración (0pt para eliminarla)
    \renewcommand{\thesubsection}{}             % Elimina el número de subsección en el índice
    
    % Subsubsecciones
    \renewcommand{\cftsubsubsecpresnum}{}       % Elimina el prefijo de numeración
    \renewcommand{\cftsubsubsecaftersnum}{}     % Elimina el sufijo de numeración
    \renewcommand{\cftsubsubsecnumwidth}{0 pt}  % Ancho de la numeración (0pt para eliminarla)
    \renewcommand{\thesubsubsection}{}          % Elimina el número de subsubsección en el índice